\chapter{Entropy--topology separation}\label{ch:separation}

The central question of this chapter: does counting histories, or counting them with their local neighbourhoods, tell an observer whether the history space has a hole? The answer has three layers. (i) \emph{In general no}: there are information-equivalent, topologically inequivalent pieces at every radius $\le2$ (\cref{sec:witnesses}). (ii) \emph{In two exactly characterised arithmetic families} the history count alone suffices, because it is injective on the comparable window (\cref{sec:window}). (iii) \emph{On the six-label census} the count alone does not suffice, but radius-two histograms do, with zero residual entropy (\cref{sec:census-entropy}); the status of that zero is carefully limited. Predictive questions are deferred to \cref{ch:predictive}.

\section{Information-equivalent, topologically inequivalent pieces}\label{sec:witnesses}
By a \emph{piece} we mean a content with a source (a regime) whose admissible set is a single component, with the data of \cref{sec:ladder}. Two pieces are \emph{information-equivalent at radius $r$} if they have the same fine accessibility signature through radius $r$: the same number of histories, the same class sizes and the same multiset of ball sizes $|B_i(w)|$ for $i\le r$, each ball considered together with its vertex's local profile.
\begin{proposition}[information-equivalent, topologically inequivalent pieces]\label{prop:witnesses}
There are pairs of arithmetic pieces with equal $N$, equal class structure (hence equal Hartley entropy $\log_2N$ and equal class entropies), agreeing in the fine signature through radius $s-1$ and differing in Betti vectors; $s$ is the \emph{separating radius}:
\begin{center}\footnotesize\setlength{\tabcolsep}{3pt}
\begin{tabular}{@{}llllcl@{}}\toprule
ID & piece A (content @ source) & $\beta(A)$ & piece B & $\beta(B)$ & $s$ and what differs\\ \midrule
W1 & $\{2,3,14,15\}@210$ & $(1,1,0,0)$ & $\{1,2,3,4\}@6$ & $(1,0,0,0)$ & $1$: $E=20$ vs $21$\\
W2 & $\{2,3,8,9\}@72$ & $(1,1,0,0)$ & $\{2,3,5,6\}@30$ & $(1,0,0,0)$ & $1$: $E=20$ vs $21$\\
W3 & $\{2,3,8,9,14\}@1008$ & $(1,0,1,0,0)$ & $\{1,2,3,5,7\}@105$ & $(1,0,0,0,0)$ & $2$: squares $66$ vs $60$\\
W4 & $\{2,3,8,9,10\}@72$ & $(1,1,0,0,0)$ & $\{1,2,3,5,9\}@18$ & $(1,0,0,0,0)$ & $2$: squares $21$ vs $19$\\
W5 & $\{2,3,9,15,21\}@189$ & $(1,1,0,0,0)$ & $\{1,2,12,13,26\}@52$ & $(1,0,0,0,0)$ & $3$: hexagons $1$ vs $2$\\
W6 & $\{2,3,8,21,26\}@273$ & $(2,1,0,0,0)$ & $\{6,8,9,14,15\}@168$ & $(2,0,0,0,0)$ & $3$: hexagons $0$ vs $1$\\ \bottomrule
\end{tabular}
\end{center}
with $N=16,16,96,40,24,18$ histories respectively. In W5 and W6 the two pieces have equal $N$, $E$ (32 and 17), square count ($7$; none) and equal fine statistics through radius $2$; their Betti numbers differ by exactly one hexagonal $2$-cell, which no statistic of radius $\le2$ sees. The Euler characteristic $\sum(-1)^if_i=\sum(-1)^i\beta_i$ holds on every row.
Consequently: histories-counting ($r=0$), one-swap ($r\le1$) and two-swap ($r\le2$) accessibility do not detect history homology on arithmetic pieces.
\status{\CompAss\ (existence statement with explicit finite witnesses; each piece is built by simulating the admissibility rule). The manuscript verifier reproduces $N$, the $f$-vectors and Betti numbers of W1, W2, W5 and W6 from the definition (\texttt{ch10\_witnesses}); W3, W4 and the separating radii are inherited and recorded \texttt{NOT\_RERUN}}{GIG-I v001, T9 (certificates W1--W6); T8}
\end{proposition}
The case W5/W6 is the sharpest: radius-two data determine $N$, $E$ and $Q$ and the identity of \cref{thm:radius2} contains no hexagon term, and these pieces differ by one hexagonal face. The combinatorial source of the difficulty is recorded in \cref{prop:petrie}: a hexagon that is the boundary of a braid face and a six-cycle that is not cannot be told apart by the graph alone without the cube-product property.

\begin{remark}[a conjecture that is open]
In the arithmetic census of distinct-label pieces the fine signature \emph{through radius $3$} determines the permutohedral Betti vector for $k\le5$ (no counterexample for $k=4$, labels $\le70$, and $k=5$, labels $\le26$); radius $2$ provably fails (W5, W6). Whether radius $3$ suffices in general is open. This is the first information-recovery conjecture of the programme; it is recorded as \cref{prob:radius3}.
\end{remark}

\section{The comparable window: where the count suffices}\label{sec:window}
Return to the anchors of \cref{ch:topology}: $\lambda_A=(2,3,4,5,8,9)$ with $m_0=288$ and $\lambda_B=(2,3,6,7,8,9)$ with $m_0=432$. The sources $m=m_0t$ with $t\in\{1,\dots,60\}$ realise only a few of the finite lattice of systems (\cref{prop:finite-lattice}: $56$ classes for $\lambda_A$, $60$ for $\lambda_B$; the sixty sources realise $8$ of them). Let $Y_{\mathrm{anc}}$ be the anchored hole state of \cref{sec:ladder}.
\begin{observation}[exact survival entropy and recovery on the window]\label{obs:anchored}
Under the uniform measure on $t=1,\dots,60$:
\begin{enumerate}[label=(\roman*)]
\item $Y_{\mathrm{anc}}=1$ exactly for $\gcd(t,30)=1$ ($16/60$; $p=4/15$) for $\lambda_A$ and for $\gcd(t,42)=1$ ($17/60$) for $\lambda_B$ (\cref{thm:anchor}); $H(Y)=-p\log_2p-(1-p)\log_2(1-p)=0.83664$ bits for $\lambda_A$ and $0.85995$ bits for $\lambda_B$.
\item The eight system classes on the window have eight \emph{distinct} history counts: $372,408,444,480,612,648,684,720$ ($\lambda_A$) and $276,312,324,360,612,648,684,720$ ($\lambda_B$) --- recomputed by the manuscript verifier from the admissibility rule (\texttt{ch10\_anchor\_A\_window}, \texttt{ch10\_anchor\_B\_window}). Therefore $H(Y\mid N)=0$: on the comparable window the history count alone recovers the hole state, and so do all finer levels. There is no pair of sources on the window with equal history count and different hole state.
\item Beyond the window, over all system classes, $\lambda_A$ has no binary separating pair at any level. $\lambda_B$ has exactly one at the level of the count: $m=648$ ($N=276$, state $\mathrm{NA}$) and $m=432$ ($N=276$, state $1$); it is separated by the ball histograms. This is not a ``filled vs survives'' pair: $648$ is not a multiple of $432$, so $Y_{\mathrm{anc}}(648)$ is undefined and counted as $0$ only by the binary convention.
\item \emph{Arrow law.} On the window, multiplying the source by a prime of the modulus ($2,3,5$ for $\lambda_A$; $2,3,7$ for $\lambda_B$) sends $Y=1\mapsto0$ at once, and $Y=0$ stays $0$: the filled state is absorbing (for divisibility $m\to mp$, not for time). The exact information jump from the surviving source is, for $\lambda_A$, $\log_2(408/372),\log_2(444/372),\log_2(612/372)=0.1333,0.2553,0.7182$ bits for $\times2,\times3,\times5$, and for $\lambda_B$ $\log_2(312/276),\log_2(324/276),\log_2(612/276)=0.1769,0.2313,1.1489$ bits for $\times2,\times3,\times7$. Multiplication beyond the cap leaves $A$, $N$ and $Y$ unchanged.
\end{enumerate}
\status{(i): \Analytic\ given \cref{thm:anchor} and \FiniteExact\ for the entropy values; (ii) \FiniteExact\ (counts re-derived); (iii), (iv): \FiniteExact\ (inherited, \texttt{NOT\_RERUN}); whether injectivity of $N$ on the window is a theorem is \textbf{open}, it is a table}{GIG-INFO v001, GI-1--GI-4, GI-7; GI-8 open}
\end{observation}
\begin{remark}[how to read this]
The injectivity of $N$ on the window is a statement about the counting law of \cref{obs:counting}: for $\lambda_A$, $N(288t)=372+36[2|t]+72[3|t]+240[5|t]$ is an \emph{injective} function of the triple of indicators on the window, because the increments $36,72,240$ have distinct subset sums. This is a numerical coincidence of the family, not a mechanism; for $\lambda_B$ the non-additive $3\cdot7$ interaction ($+24$) does not destroy injectivity. A general-grammar version of the observation hierarchy for non-integer systems is open.
\end{remark}

\section{The census: the count does not suffice, radius two does}\label{sec:census-entropy}
Let $\mathcal D$ be the six-label census: $8{,}008$ contents with labels in $\{1,\dots,16\}$, $346{,}859$ regimes (distinct nonempty admissibility sets). Let $Y=\mathbf 1[\beta_2(K_A;\F)>0]$.
\begin{observation}[the count does not determine $Y$ on the census]\label{obs:equal-N}
The recomputation of every Betti vector with a second engine gives $158$ regimes with $\beta_2>0$ ($157$ with $\beta_2=1$, one with $\beta_2=3$), identical as a set to the frozen list of \cref{obs:census-H2}. The $158$ positives take $59$ distinct values of $N$ (ranging over $1\ldots720$). There are $326{,}912$ (positive, negative) pairs with equal $N$; $72$ of them lie in the same content, and $51$ positives have a same-content negative partner. Every one of the $158$ positives has at least one negative with equal $N$. Hence $H(Y\mid N)>0$ under each of the ensembles below.
\status{\FiniteExact\ (inherited, \texttt{NOT\_RERUN})}{GIG-INFO v002, GI2-1, GI2-2}
\end{observation}
\begin{example}[a certified pair with equal count]\label{ex:pair}
Content $(2,3,4,5,8,15)$ at sources $480$ and $800$: both have $N=396$ (re-derived by simulating divisibility), but $f$-vectors $(396,900,669,173,10,0)$ and $(396,882,636,156,7,0)$ and Betti vectors $(1,0,1,0,0,0)$ and $(1,0,0,0,0,0)$; the positive carries an explicit $2$-cycle and $2$-cocycle with $\varphi(z)=1$ (\cref{lem:certificate}), the negative has $\beta_2=0$ by two independent GF(2) rank routes. The smallest exhibit with equal radius-one data is content $(2,4,5,9,14,15)$ at $2520$ ($\beta_2=1$) and $3780$ ($\beta_2=0$): both have $N=540$, equal $h_1$, equal square count $Q=738$, equal $E=1278$, equal $S_{d2}=12212$ and equal $S_B=9800$, but $S_{BB}=180420$ vs $180408$ (different $h_2$); their $f$-vectors differ in $f_3$ only ($291$ vs $292$). They share $468$ of $540$ arrangements.
\status{\FiniteExact; $N$, $f$-vectors, Betti numbers, $S_{BB}$ and $h_1,h_2$ re-derived by the manuscript's figure generator (\texttt{figures/fig5\_data.json}); certificates inherited}{GIG-INFO v002, GI2-6; v003, R3}
\end{example}
\begin{figure}[ht]\centering\includegraphics[width=\textwidth]{figures/fig5_equal_N_different_topology.pdf}
\caption{Equal history count, different topology. (a) $f$-vectors of the pair of \cref{ex:pair} with $N=396$: the surviving hole (red) has more cells than the filled complex (blue) in every dimension $\ge1$. (b,c) Radius-two ball-size histograms of the pair $2520$ vs $3780$ at $N=540$ with equal radius-one data: the histograms differ in a few vertices only. All values are recomputed exactly.}\label{fig:equal-N}
\end{figure}
\begin{theorem}[no function of the radius-one data and the square-incidence data determines $Y$]\label{thm:no-radius1}
There is no function of $(N,h_1,h_q)$ --- hence none of $(N,E,Q,S_{d2},S_B)$ --- whose value determines $Y$ on $\mathcal D$. Indeed there are $165$ (positive, negative) pairs with the same $(N,h_1,h_q)$, involving $22$ positives, and all have different $h_2$; the pair of \cref{ex:pair} is one.
\status{\FiniteExact\ (witness), \CompAss: the statement is the existence of the explicit pair above; the count $165/22$ is inherited, \texttt{NOT\_RERUN}}{GIG-INFO v003, R3}
\end{theorem}
\begin{table}[ht]\centering\footnotesize\setlength{\tabcolsep}{3pt}
\begin{tabular}{@{}lrrrrrrrr@{}}\toprule
ensemble & size & $H(Y)$ & $H(Y|N)$ & $H(Y|N,f_1)$ & $H(Y|N,\chi)$ & $H(Y|N,h_1)$ & $H(Y|f)$ & $H(Y|N,h_2)$\\ \midrule
$\mathcal E_{\mathrm{all}}$ & $346{,}859$ & $.005713$ & $.004467$ & $.002172$ & $.000109$ & $.000619$ & $0$ & $0$\\
$\mathcal E_{\mathrm{band}}$ & $93{,}018$ & $.018078$ & $.016659$ & $.008100$ & $.000406$ & $.002309$ & $0$ & $0$\\
$\mathcal E_{\mathrm{rank}}$ & $219{,}146$ & $.008231$ & $.006276$ & $.003140$ & $.000168$ & $.000803$ & $0$ & $0$\\ \bottomrule
\end{tabular}
\caption{Exact residual entropy (bits) of $Y=\mathbf 1[\beta_2>0]$ given observations, for three uniform ensembles on the census (full 30-digit values are in the source records). $\mathcal E_{\mathrm{band}}$: regimes whose $N$ is attained by a positive regime. $\mathcal E_{\mathrm{rank}}$: one regime per distinct \emph{rank-indexed} family mask ($219{,}146$; the source register's earlier name ``distinct arrangement sets'' was a mislabel, corrected in the third information release). $f$ denotes the full $f$-vector. Plug-in values, in-sample, see \cref{ch:predictive}.}\label{tab:residual}
\end{table}
\begin{observation}[radius-two data recover $Y$ on the census; the zero is a collision statement]\label{obs:residual}
On $\mathcal D$ the observation $O_h=(N,h_1,h_2)$ has zero mixed classes in all three ensembles ($8{,}208$ classes in $\mathcal E_{\mathrm{all}}$, none mixed), so $H(Y\mid O_h)=H(Y\mid O_G)=H(Y\mid O_A)=0$ (\cref{thm:obs-ladder}). The separation $H(Y\mid N)>0$ therefore does \emph{not} persist once radius-two accessibility is added. Of partial observables, $(N,h_2)$ and the full $f$-vector already give $0$, while $(N,h_1)$, $(N,f_1)$, $(N,\chi)$ leave $66$, $151$, $8$ positives sharing their signature with a negative.
These are empirical facts about $\mathcal D$ ($\le16$ labels, six distinct labels). They are not theorems, they do not extend to other label ranges or sizes, and their arithmetic origin is \textbf{open} (\cref{prob:arith-origin}). Moreover the zero residual is a statement about \emph{collisions} of signature classes in a finite census; it is not a statement about information that generalises (\cref{ch:predictive}).
\status{\FiniteExact\ (inherited, \texttt{NOT\_RERUN}); \Bounded\ as to scope}{GIG-INFO v002, GI2-3, GI2-5}
\end{observation}
\begin{remark}[Hartley versus Shannon, summarised]
Two quantities are in play. The \emph{Hartley} ambiguity $\log_2N$ of a regime measures the multiplicity of its histories. The \emph{Shannon} residual $H(Y\mid\cdot)$ measures what a given observation leaves unknown about topology. They are unrelated in general: \cref{prop:witnesses} and \cref{obs:equal-N} exhibit equal Hartley and different topology, and the zero of \cref{obs:residual} exhibits zero residual with positive Hartley. The unifying fact is the ladder of \cref{thm:obs-ladder}: every observation is a function of $A$, and topology is a function of $A$ that is \emph{not} a function of the cheap observations.
\end{remark}
